\documentclass[10pt,a4paper]{amsart}
\usepackage[margin=19mm]{geometry}
\usepackage{amsmath,amssymb,mathtools}
\usepackage[scr=rsfs,cal=euler]{mathalfa}
\usepackage{xfrac}
\usepackage[unicode,hidelinks,bookmarks=false]{hyperref}
\newcommand{\ie}{i.e.\ }
\newcommand{\cf}{cf.\ }
\newcommand{\resp}{respectively\ }
\newcommand{\Subsection}{\subsection}
\providecommand{\nobreakdash}{\nobreak\hbox{-}}
\setlength{\emergencystretch}{2em}
\multlinegap0pt
\usepackage{bm}
\usepackage{bbm}
\usepackage{dsfont}
\usepackage{xspace}
\usepackage{cancel}
\usepackage{mathtools}
\usepackage{amsthm}
\makeatletter
\@ifundefined{theorem}{\newtheorem{theorem}{Theorem}}{}
\@ifundefined{lemma}{\newtheorem{lemma}[theorem]{Lemma}}{}
\makeatother
\newcommand{\Prob}{\mathbb{P}}
\title[Limit Laws for Consensus Protocols on the Complete Graph]{Limit Laws for Consensus Protocols on the Complete Graph}
\author{Julian Becker, Konstantinos Panagiotou}
\date{Source: arXiv:2605.19131v2 (CC BY 4.0)}
\begin{document}
\maketitle

\noindent\emph{Source and licence.} Limit Laws for Consensus Protocols on the Complete Graph. Version arXiv:2605.19131v2 (CC BY 4.0), distributed under the Creative Commons Attribution 4.0 International licence, as recorded in the arXiv licence metadata for this exact version and shown on the abstract page https://arxiv.org/abs/2605.19131v2. Full rights notice: licenses/arxiv-2605.19131-CC-BY-4.0.txt.\par\smallskip

\noindent\textit{Editorial scope.} Counted unit: the Lemma whose source label is \texttt{lem: ending phase concentration} in the article (source0.tex lines 1018--1024, source label \texttt{lem: ending phase concentration}), delivered together with the article's own complete original proof of it (source0.tex lines 1163--1263, one \texttt{proof} environment of 101 lines). No mathematics is written, repaired or re-derived: statement and proof are the article's own text line for line. The proof is detached from the statement in the source: the article states the result at lines 1018--1024 and prints its proof at lines 1163--1263, and the proof environment's own header names the statement, so the association is the article's own. Required context retained uncounted: lem: threshold sqrt(n) (lines 1030--1039); lem: aux lem 2.5 (lines 1102--1111). Every definition, display and label of those lines is retained unchanged. Omitted from this package and not redistributed: 1--1017, 1025--1029, 1040--1101, 1112--1162, 1264--1695. Nothing inside a retained excerpt is omitted or altered: every retained body is delivered exactly as the article's lines are written, so no omission receipt of any kind accompanies this package. Citations: the retained excerpts carry no citation command at all, so no bibliography entry of the article is transcribed and no reference is redistributed. Reference adaptation: none was needed and none was made; every reference inside the delivered unit resolves inside it, so no unresolved reference or citation is produced. Printed slips kept literal: no printed word, symbol, hypothesis or number of the counted statement or of its proof was altered. Layout and class: the article's own class is \texttt{article} with packages including graphicx, float, caption, subcaption, mathtools, mathrsfs, verbatim, amsmath, amssymb, amsthm, bbm, xcolor, hyperref, inputenc; the wrapper is the lane's pinned \texttt{amsart} editorial wrapper at 10pt on A4, which restates the theorem-like environments the retained text needs and adds the article's own macro definitions that the retained text needs (\texttt{\textbackslash Prob}), with extra wrapper packages (bm, bbm, dsfont, xspace, cancel, mathtools). The delivered numbering is the unit's own. Assets: no retained span contains a figure environment, a graphics-inclusion command, a PDF-page inclusion, a table or a TikZ picture; no image file is copied into this package, the package declares no asset dependency and no image file is redistributed with it. Licence: version arXiv:2605.19131v2 of this article is distributed under the Creative Commons Attribution 4.0 International licence, as recorded in the arXiv licence metadata for this exact version and shown on the abstract page https://arxiv.org/abs/2605.19131v2. A full rights notice is delivered at licenses/arxiv-2605.19131-CC-BY-4.0.txt. Characters: every retained excerpt is pure ASCII, so the delivered engine, XeLaTeX with its default Latin Modern text font, reports no missing character.
\begin{lemma}\label{lem: threshold sqrt(n)}
    Let $t, t' \in \mathbb{N}$ such that $Y_t \geq n^{1-1/m}\ln n$ and $Y_{t'} \leq n^{1-1/m}/\ln n$. Then
    \[
    \Prob(Y_{t+1} = 0 \mid Y_t) = o(1)
    \quad 
    \text{and}
    \quad
    \Prob(Y_{t'+1} > 0\mid Y_{t'}) = o(1).
    \]
\end{lemma}

\begin{lemma}\label{lem: aux lem 2.5}
Let $0 < \delta < 1/4$. Then, the event $\bigcap_{0 \le t \le 2\log_m\ln n}C_t$ occurs whp, where
\[
    C_t := \left \{\left|Y_{t+1} - \mathbb{E}[Y_{t+1}\mid Y_t] \right| \leq  \mathbb{E}[Y_{t+1}\mid Y_t]^{1/2+\delta} + n^\delta \right \}
    \quad
    \text{and}
    \quad
    \mathbb{E}[Y_{t+1}\mid Y_t] = f(Y_t/n)n.
\]
\end{lemma}

\begin{lemma}
\label{lem: ending phase concentration}
Let $t^\star\in\mathbb{N}$ and $X_{t^\star}=n-Y_{t^\star} = x n$ for $1/2<x<1$.  Then, whp
\[
    \bigcap_{0 \leq t \leq 2 \log_m\ln n} \left\{\left|Y_{t^\star+t} - f^{(t)}(Y_{t^\star}/n)n \right|\leq \max\left\{Y_{t^\star+t}n^{-1/20}, n^{1/4}\right\} \right\}.
\]
\end{lemma}

\begin{proof}[Proof of Lemma \ref{lem: ending phase concentration}]
Since $(X_t)_{t\in \mathbb{N}_0}$ is markovian we can assume without loss of generality that $t^\star = 0$. In particular, $X_0 = xn$ and $Y_0 = (1-x)n$.
We abbreviate
\[
    \Delta_{t} = \left|Y_{t} - f^{(t)}(Y_0/n)n \right|,
    \quad\text{for } 0 \le t \le 2 \log_m\ln n.
\]
For later reference, since $f$ is of majority-type, we have $f(x)=\beta x^m + O(x^{m+1})$ as $x\to 0$. Moreover, the first $m-1$ derivatives of $f$ in zero are zero, since the forward difference representation shows
\begin{align*}
    \frac{d^s}{dx^s} f (0) &= \lim_{h\searrow 0} \frac{\sum_{j=0}^{s} (-1)^{s-j}\binom{s}{j}f(jh)}{h^s} = \lim_{h\searrow 0} \sum_{j=1}^{s} (-1)^{s-j} j^s \binom{s}{j} \frac{f(jh)}{(jh)^s} = 0 
    \quad\text{for }
    1\leq s <m.
\end{align*}
It follows by Taylor's theorem that
\begin{align}\label{equ: formula derivative}
    f'(x)=m\beta x^{m-1} + O(x^{m}).
\end{align}
Hence, there exist $c_u,c_\ell>0$ such that for $x\in[0,1]$,
\begin{align}\label{inequ: f upper lower x^m}
    c_\ell x^m \leq f(x) \leq c_ux^m
    \quad\text{and}\quad
    f'(x)\leq c_ux^{m-1}.
\end{align}
%{\color{red}
Moreover, if $Y_{t'} < n^{1/4}$ for some $t' \in \mathbb{N}$, then $Y_{t'+1} =0$ whp by Lemma \ref{lem: threshold sqrt(n)}. Hence, let $\overline{T} \in \mathbb{N}_0 \cup \{\infty\}$ be such that, whp,
\begin{align}\label{consequ of Lem 2.8}
    Y_{t} \geq n^{1/4} ~ \text{for all} ~  t\leq \overline{T}
    \quad \text{and} \quad
    Y_{t} < n^{1/4} ~ \text{for all} ~  t > \overline{T}.
\end{align}
%}
In order to prove the lemma, we condition
%{\color{red}
on the event \eqref{consequ of Lem 2.8} and additionally
%}
on the event in Lemma~\ref{lem: aux lem 2.5}, where, with foresight, $\delta=1/14$. We claim for $n$ large enough and $\lambda:=2^{m+1}c_u/c_\ell$ that
%{\color{red}
%\[
\begin{equation}
\label{eq:boundD_t}
    \Delta_t \leq \lambda^tY_t/n^{1/14},
    ~~\text{for}~~
    t\leq \min\{\overline{T}, 2\log_m\ln n\}
    \quad\text{and}\quad
    \Delta_{t} \leq n^{1/4}
    ~~\text{for}~~
    \overline{T} < t \le 2\log_m\ln n .
\end{equation}
This implies the statement of the lemma, since $\lambda^t = n^{o(1)}$ for $t \leq 2 \log_m \ln n$.
Thus, it remains to prove~\eqref{eq:boundD_t}. We proceed by induction over $t \in \mathbb{N}_0$. Note that $\Delta_{0}=0$, establishing the base case. For the induction step, we assume that the claim holds for some $t\in \mathbb{N}_0$.
Recall that $\mathbb{E}[Y_{t+1}\mid Y_t] = f(Y_t/n)n$. Since $Y_t$ satisfies Lemma~\ref{lem: aux lem 2.5} with $\delta = 1/14$, applying the triangle inequality implies
\begin{equation}\label{inequ: first bound delta}
\begin{aligned}
    \Delta_{t+1} &\leq \left|Y_{t+1}-f(Y_t/n)n \right| +\left|f(Y_t/n)n -  f^{(t+1)}(Y_0/n)n \right| \\
    &\leq (f(Y_t/n)n)^{4/7} + n^{1/14} + n\left|f(Y_t/n) -  f^{(t+1)}(Y_0/n) \right|.
\end{aligned}
\end{equation}
By definition $Y_t/n = f^{(t)}(Y_0/n) \pm \Delta_{t}/n$. Hence, using the mean value theorem, there exists $\xi$ between $Y_t/n$ and $f^{(t)}(Y_0/n)$, so that $0\leq \xi \leq Y_t/n + \Delta_t/n$, and such that
\begin{align*}
\label{inequ: second summand1}
    n \left|f(Y_t/n) -  f^{(t+1)}(Y_0/n) \right|
    = f'(\xi)\Delta_t
    \stackrel{\eqref{inequ: f upper lower x^m}}{\leq}c_u\, (Y_t/n +\Delta_t/n)^{m-1}  \Delta_t.
\end{align*}
This implies that 
\begin{equation}
\label{inequ:2ndboundD_t}
    \Delta_{t+1}
    \le
    (f(Y_t/n)n)^{4/7} + n^{1/14} + c_u\, (Y_t/n +\Delta_t/n)^{m-1}  \Delta_t.
\end{equation}
In what follows we distinguish whether $t > \overline{T}$ or $t \le \overline{T}$.
Suppose $t > \overline{T}$. 
By~\eqref{consequ of Lem 2.8} we have $Y_t\leq n^{1/4}$ and the induction hypothesis yields $\Delta_t\leq n^{1/4}$, too. 
From~\eqref{inequ:2ndboundD_t} we obtain that
    \begin{align*}
    \Delta_{t+1}
    %&\leq (c_u(Y_{t^\star+t}/n)^mn))^{4/7} + n^{1/14} + c_u\, (Y_{t^\star+t}/n +\Delta_{t}/n)^{m-1}  \Delta_{t}\\
    &\leq c_u^{4/7}n^{(4-3m)/7}+n^{1/14}+c_u\,2^{m-1}n^{(4-3m)/4},
    \end{align*}
    which is for sufficiently large $n$ and independently of $t$ at most $n^{1/4}$, as required.
    It remains to consider the case $t\leq \overline{T}$.
Recall that $\lambda=2^{m+1}c_u/c_\ell$. We proceed from~\eqref{inequ:2ndboundD_t}, where we first note that since $t\leq2\log_m\ln n$, the induction hypothesis implies $\Delta_t\leq Y_t$.
Applying $x^m\leq f(x)/c_\ell$ from~\eqref{inequ: f upper lower x^m} and using again the induction hypothesis, this time for the second occurrence of $\Delta_t$, we obtain
    \begin{equation}\label{inequ: second bound delta}
        \Delta_{t+1} \leq (f(Y_t/n)n)^{4/7} + n^{1/14} + \frac{\lambda^{t+1}}{4}  n^{13/14}f(Y_t/n).
    \end{equation}
Since $Y_t$ satisfies Lemma \ref{lem: aux lem 2.5} with $\delta=1/14$. we obtain
$f(Y_t/n)n \leq Y_{t+1}+\left(f(Y_t/n)n\right)^{4/7}+n^{1/14}$.
Note that this guarantees that $f(Y_t/n)n \le 2(Y_{t+1} + n^{1/14})$ whenever $n$ is sufficiently large; to see this, note that if $f(Y_t/n)n$ is larger than some specific constant, then $(f(Y_t/n)n)^{4/7} \le f(Y_t/n)n/2$. Thus,~\eqref{inequ: second bound delta} asserts that
\begin{equation}
\label{eq:Dt+1aux}
    \Delta_{t+1} \leq \left(2Y_{t+1} + n^{1/14}\right)^{4/7} + n^{1/14} + \frac{\lambda^{t+1}}{2} (Y_{t+1} + n^{1/14})\, n^{-1/14}.
\end{equation}
If $Y_{t+1} \ge n^{1/4}$, then this readily implies 
\[
    \Delta_{t+1}
    \le {\lambda^{t+1}}Y_{t+1}n^{-1/14} \left(\frac12 + o(1)\right)
\]
and~\eqref{eq:boundD_t} is established for large $n$. On the other hand, if $Y_{t+1} < n^{1/4}$, then the right hand side in~\eqref{eq:Dt+1aux} is already $o(n^{1/4})$, and~\eqref{eq:boundD_t} is established for large $n$ as well.
\end{proof}

\end{document}
